\documentclass[12pt,a4paper]{article}
\usepackage[T1]{fontenc}
\usepackage{amsmath}
\usepackage{amssymb}
\usepackage{amsthm}
\usepackage{geometry}
\usepackage{xcolor}
\usepackage{fancyhdr}

\geometry{top=2cm, bottom=2cm, left=2cm, right=2cm}

% Definizione stili per teoremi
\theoremstyle{definition}
\newtheorem{definition}{Definizione}[section]
\newtheorem{proposition}[definition]{Proposizione}
\newtheorem{theorem}[definition]{Teorema}

% Colori
\definecolor{darkblue}{RGB}{31, 78, 120}
\definecolor{lightblue}{RGB}{46, 92, 138}

% Intestazione
\pagestyle{fancy}
\lhead{Algorithmic Game Theory - Lecture 22}
\rhead{}
\cfoot{\thepage}

\title{\textcolor{darkblue}{\Large Algorithmic Game Theory \\ Lecture 22 \\ Coalitional Structures and Core Approximations}}
\author{Laurea Magistrale in Computer Science \\ 2025/26}
\date{}

\begin{document}

\maketitle

\section*{Slide 1: Introduzione}

Questo documento contiene gli appunti della Lezione 22 del corso di Algorithmic Game Theory.

\textbf{Argomento principale:} The Core for Coalitional Structures and Core Approximations

\newpage

\section{The Core for Coalitional Structures}

\subsection*{Strutture Coalitive}

Supponiamo che vengano formate diverse coalizioni disgiunte: $\mathcal{P} \in \mathcal{P}(N)$.

Una struttura coalitiva e una collezione di coalizioni disgiunte che partizionano l'insieme dei giocatori.

\subsection*{Imputazioni per Strutture Coalitive}

\begin{definition}[Imputazione per struttura coalitiva]
Un vettore $x \in \mathbb{R}^n$ e un'imputazione per $(N, v)$ con rispetto a $\mathcal{P} \in \mathcal{P}(N)$ se e:
\begin{itemize}
    \item \textbf{Socialmente razionale (efficiente)}: $x(S) = v(S)$ per tutti i $S \in \mathcal{P}$
    \item \textbf{Individualmente razionale}: $x_i \geq v(i)$ per ogni giocatore $i \in N$
\end{itemize}
\end{definition}

L'insieme delle imputazioni e denotato da:
\[
X^n(P, v) = \{\text{imputazioni}\}
\]

\subsection*{Il Core di un Gioco}

\begin{definition}[Core di un gioco]
\[
C(N, v, \mathcal{P}) = \left\{ x \in X^n(\mathcal{P}, v) : x(S) \geq v(S) \text{ per ogni } S \subseteq N \right\}
\]

Imputazioni che sono razionali per ogni possibile coalizione.
\end{definition}

\newpage

\section{Superadditivity Drives Likely Coalitional Structures}

\subsection*{Superadditività e Coalizioni}

Sia $(N, v^*)$ il \textbf{superadditivo cover} di $(N, v)$, definito da:
\[
v^*(S) = \max_{T \in \mathcal{P}(S)} \sum_{T \in \mathcal{T}} v(T)
\]

\begin{theorem}[Strutture coalitive in giochi superadditivi]
\[
C(N, v, \mathcal{P}) = C(N, v^*) \cap X(\mathcal{P}, v)
\]
\end{theorem}

\subsection*{Quando il Core e Non-Vuoto}

Per una partizione $\mathcal{P} \in \mathcal{P}(N)$:

\begin{itemize}
    \item Se $v^*(N) = \sum_{T \in \mathcal{P}} v(T)$, allora $C(N, v, \mathcal{P}) = C(N, v^*)$

    \item Se $v^*(N) > \sum_{T \in \mathcal{P}} v(T)$, allora $C(N, v, \mathcal{P}) = \emptyset$
\end{itemize}

Quando il grande gioco offre piu valore che le coalizioni separate, il core della struttura coalitiva non e vuoto.

\newpage

\section{Market Games with Agreement Costs}

\subsection*{Configurazione del Gioco}

Consideriamo un'economia con costi di accordo:
\begin{itemize}
    \item $N = P \cup B$ con $P = \{\text{produttori di vino}\}$ e $B = \{\text{produttori di bottiglie}\}$
    \item $l_i$: litri di vino prodotto da $i \in P$
    \item $b_i$: bottiglie prodotto da $i \in B$
    \item $c$: costo di transazione per bottiglia prodotta
\end{itemize}

\subsection*{Funzione Caratteristica con Costi}

\[
v(S) = \min \left\{ \sum_{i \in S \cap P} l_i, \sum_{j \in S \cap B} b_j \right\} - c |S \cap P| \cdot |S \cap B|
\]

Il termine $-c |S \cap P| \cdot |S \cap B|$ rappresenta il costo totale di negoziazione tra produttori di vino e produttori di bottiglie nella coalizione $S$.

\subsection*{Proprieta Fondamentali}

Quando il costo di accordo e sufficientemente alto:
\begin{itemize}
    \item $C(N, v) = \emptyset$ (nessuna soluzione stabile)
    \item Il core e vuoto perche il costo di coordinarsi supera i benefici
\end{itemize}

Se $v^*(N) = \sum_{T \in \mathcal{P}} v(T)$ (le coalizioni separate non hanno costi):
\begin{itemize}
    \item $\mathcal{P} = \{\{1, 3\}, \{2, 4\}\}$ (le partizioni ottimali)
    \item Oppure $\mathcal{P} = \{\{2, 3, 5\}, \{1, 4\}\}$ (altre configurazioni equilibrate)
\end{itemize}

\newpage

\section{Wishing to Be Back to a Unique Coalition}

\subsection*{Setup Classico}

Torniamo al problema classico del core:
\begin{itemize}
    \item $N = \{1, \ldots, n\}$: $n$ giocatori che cercano un accordo vincolante
    \item Coalizione: $S \subseteq N$
    \item Funzione caratteristica: $v : \mathcal{P}(N) \to \mathbb{R}$ (con $v(\emptyset) = 0$)
    \item $v(S)$ = worth (valore) della coalizione $S$
\end{itemize}

\subsection*{La Grande Domanda}

Supponiamo che la coalizione $N$ di tutti i giocatori si formi.

\textcolor{orange}{\textbf{Come splitzeranno i giocatori il valore $v(N)$ tra loro?}}

\subsection*{Preimputazioni}

\begin{definition}[Preimputazione]
Un vettore $x \in \mathbb{R}^n$ e una \textbf{preimputazione} per $(N, v)$ se e socialmente razionale, cioe $x(N) = v(N)$.
\end{definition}

L'insieme delle preimputazioni e:
\[
X^n(N, v) = \{\text{preimputazioni}\}
\]

Questo e un iperpiano e e sempre non-vuoto.

\textcolor{orange}{\textbf{Quali preimputazioni sono ragionevoli/accettabili?}}

\subsection*{Il Core di un Gioco}

\begin{definition}[Core]
\[
C(N, v) = \left\{ x \in X^n(N, v) : x(S) \geq v(S) \text{ per ogni } S \subseteq N \right\}
\]
\end{definition}

\newpage

\section{Approximating the Core: Quasi Cores}

\subsection*{Il Strong Epsilon-Core}

\begin{definition}[Strong epsilon-core]
Secondo Shapley-Shubik 1966:
\[
C_\varepsilon(N, v) = \left\{ x \in X^n(N, v) : x(S) \geq v(S) - \varepsilon \text{ per ogni } S \subseteq N \right\}
\]

dove $\varepsilon \in \mathbb{R}$.
\end{definition}

\subsection*{Proprieta dell'Epsilon-Core}

\begin{itemize}
    \item $C_0(N, v) = C(N, v)$ (il core classico)

    \item Per $\varepsilon_1 \leq \varepsilon_2$: $C_{\varepsilon_1}(N, v) \subseteq C_{\varepsilon_2}(N, v)$

    \item $C_\varepsilon(N, v) \neq \emptyset$ quando $\varepsilon$ e abbastanza grande

    \item $C_\varepsilon(N, v) = \emptyset$ quando $\varepsilon$ e abbastanza piccolo (eventualmente negativo)
\end{itemize}

\newpage

\section{The Weak Epsilon-Core}

\begin{definition}[Weak epsilon-core]
Secondo Shapley-Shubik 1966:
\[
\bar{C}_\varepsilon(N, v) = \left\{ x \in X^n(N, v) : x(S) \geq v(S) - |S| \cdot \varepsilon \right\}
\]
\end{definition}

\subsection*{Proprieta del Weak Epsilon-Core}

\begin{itemize}
    \item Per $\varepsilon \geq 0$: $\bar{C}_\varepsilon(N, v) \subseteq C_\varepsilon(N, v) \subseteq C_{\varepsilon}(N, v)$

    \item La relazione e invertita per $\varepsilon \leq 0$
\end{itemize}

\newpage

\section{The Least Core}

\subsection*{Definizione del Least Core}

Sia $I(N, v) = \{\varepsilon \in \mathbb{R} : C_\varepsilon(N, v) \neq \emptyset\}$ l'insieme dei valori per cui il core non e vuoto.

\begin{definition}[Least core]
Secondo Maschler-Peleg-Shapley 1979:
\[
LC(N, v) = \bigcap_{\varepsilon \in I(N, v)} C_\varepsilon(N, v)
\]
\]
\end{definition}

\subsection*{Proprieta del Least Core}

\begin{itemize}
    \item $LC(N, v) \neq \emptyset$ (sempre non-vuoto)

    \item $C(N, v) \neq \emptyset$ se e solo se $\bar{\varepsilon} \leq 0$, dove $\bar{\varepsilon}$ e il valore minimo

    \item Equivalentemente: $LC(N, v) \subseteq C(N, v)$ se e solo se il core e non-vuoto
\end{itemize}

\subsection*{Caratterizzazione Alternativa}

\begin{proposition}
\[
LC(N, v) = \arg \min \left\{ \max\{v(S) - x(S) : S \subseteq N\} : x \in X^n(N, v) \right\}
\]
\end{proposition}

Il least core minimizza il massimo disagio di qualsiasi coalizione.

\newpage

\section{Market Games with Agreement Costs and Quasi Cores}

\subsection*{Analisi con Costi di Accordo}

Riprendiamo l'economia con costi di accordo:
\[
v(S) = \min \left\{ \sum_{i \in S \cap P} l_i, \sum_{j \in S \cap B} b_j \right\} - c |S \cap P| \cdot |S \cap B|
\]

\subsection*{Quando il Core e Non-Vuoto}

\begin{itemize}
    \item Quando $c$ e grande (molti costi di accordo): $C(N, v) = \emptyset$

    \item Quando $c = 1$: il valore minimo $\bar{\varepsilon} = 1.5$

    \item Il least core: $LC(N, v) = \{(t, 0.25 - t, 3.15 - t) : 0 \leq t \leq 1\}$
\end{itemize}

Il least core fornisce un modo per trovare allocazioni ragionevoli anche quando il core e vuoto.

\newpage

\section{The Least Core and Measure of Dissatisfaction}

\subsection*{Il Least Core Revisitato}

\[
LC(N, v) = \arg \min \left\{ \max\{v(S) - x(S) : S \subseteq N\} : x \in X^n(N, v) \right\}
\]

\subsection*{Eccesso di una Coalizione}

\begin{definition}[Eccesso di una coalizione]
L'eccesso di una coalizione $S \subseteq N$ al vettore $x \in \mathbb{R}^n$ e:
\[
e(S, x) = v(S) - x(S)
\]
\end{definition}

Questo rappresenta il disagio della coalizione $S$.

\subsection*{Interpretazione dell'Eccesso}

\begin{itemize}
    \item $e(S, x) \leq 0$: la coalizione e soddisfatta

    \item $e(S, x) > 0$: il valore fornisce una misura di insoddisfazione
\end{itemize}

\subsection*{Il Core via Eccessi}

\begin{definition}[Core via eccessi]
\[
C_\varepsilon(N, v) = \left\{ x \in X^n(N, v) : e(S, x) \leq \varepsilon \text{ per ogni } S \subseteq N \right\}
\]
\end{definition}

Il least core fornisce il massimo livello di insoddisfazione minimo possibile.

\newpage

\section{Iterate the Minimization of (Maximum) Dissatisfaction}

\subsection*{Come Un Arbitro Dovrebbe Assegnare i Payoff}

\textcolor{orange}{\textbf{Come dovrebbe un arbitro assegnare i payoff?}}

\textbf{(La selezione di una suddivisione unica $x$ e richiesta)}

\subsection*{Vettore Ordinato di Insoddisfazione}

\begin{definition}[Ordered vector of dissatisfaction]
Un vettore ordinato di insoddisfazione a $x \in \mathbb{R}^n$ e definito come:
\[
\theta(x) = (e(S_1, x), e(S_2, x), \ldots, e(S_{2^n}, x))
\]

dove $e(S_i, x) \geq e(S_{i+1}, x)$ per $i = 1, \ldots, 2^n - 1$.
\end{definition}

L'insieme di coalizioni con tutti i loro eccessi (escluso il grande gioco).

\subsection*{L'Idea Fondamentale}

Vogliamo \textbf{minimizzare l'insoddisfazione lessicograficamente}.

\subsection*{Ordine Lessicografico}

\begin{definition}[Lexicographic order]
Dati $a, b \in \mathbb{R}^m$, diciamo $a \geq_{\text{lex}} b$ se:
\begin{itemize}
    \item Esiste $k$ tale che $a_i = b_i$ per ogni $i < k$ e $a_k > b_k$

    \item Oppure $a = b$
\end{itemize}
\end{definition}

Un ordine totale (riflessiva, antisimmetrica, transitiva e convessa).

L'ordine non e continuo: $(a' >_{\text{lex}} a)$ puo comportare $(a' >_{\text{lex}} b)$ per $a' \approx a$ ma $b \neq a$.

\end{document}
